\section*{अध्याय \ref{ch:g5:series-parallel-circuits} --- श्रेणी और समांतर परिपथ}

\begin{solution}{exo:g5:series-parallel-circuits:1}
\omterm{def:g5:series-parallel-circuits:parallel}{शाखा} उन अलग-अलग सड़कों में से एक है जिनमें परिपथ फूटता है; \omterm{def:g5:series-parallel-circuits:parallel}{संधि} वह
बिंदु है जहाँ सड़कें फूटती या फिर से मिलती हैं। \omterm{def:g5:series-parallel-circuits:parallel}{समांतर में} लगे दो
\omterm{def:g3:first-electric-circuit:bulb}{बल्ब} \omterm{def:g3:first-electric-circuit:bulb}{बल्ब} वाली दो \omterm{def:g5:series-parallel-circuits:parallel}{शाखाओं} पर सवार होते हैं।
\end{solution}

\begin{solution}{exo:g5:series-parallel-circuits:2}
स्वतंत्रता (दरार सिर्फ़ अपनी \omterm{def:g5:series-parallel-circuits:parallel}{शाखा} को रोकती है); पूरी चमक (हर \omterm{def:g3:first-electric-circuit:bulb}{बल्ब}
ऐसे चमकता है मानो अकेला हो, और बदले में \omterm{def:g3:first-electric-circuit:battery}{बैटरी} तेज़ी से चुकती है);
हुक्म की चौकियाँ (\omterm{def:g5:series-parallel-circuits:parallel}{शाखा} का \omterm{ex:g3:first-electric-circuit:switch}{स्विच} अपनी \omterm{def:g5:series-parallel-circuits:parallel}{शाखा} पर, और फूट से पहले लगा
\omterm{ex:g3:first-electric-circuit:switch}{स्विच} सब पर राज करता है)।
\end{solution}

\begin{solution}{exo:g5:series-parallel-circuits:3}
\omterm{def:g3:first-electric-circuit:bulb}{बल्ब} 2 बिना रुके जलता रहता है --- उसकी अपनी सड़क पूरी है। \omterm{def:g4:series-circuits:series}{श्रेणी में}
वह बुझ गया था, क्योंकि वहाँ दोनों \omterm{def:g3:first-electric-circuit:bulb}{बल्ब} एक ही अकेली सड़क साझा करते थे,
और कोई भी दरार उसे दोनों के लिए तोड़ देती थी।
\end{solution}

\begin{solution}{exo:g5:series-parallel-circuits:4}
जोड़ी ख (समांतर) में हर \omterm{def:g3:first-electric-circuit:bulb}{बल्ब} ज़्यादा चमकता है --- हर एक पूरी चमक पर।
जोड़ी ख ही अपनी \omterm{def:g3:first-electric-circuit:battery}{बैटरी} ज़्यादा तेज़ी से चुकाती है: वह एक साथ दो पूरी
\omterm{def:g5:series-parallel-circuits:parallel}{शाखाओं} को खिलाती है।
\end{solution}

\begin{solution}{exo:g5:series-parallel-circuits:5}
दोनों बंद: दोनों \omterm{def:g3:first-electric-circuit:bulb}{बल्ब} जलते हैं। मुख्य खुला: दोनों अँधेरे --- वह साझा
सड़क पर राज करता है। मुख्य बंद और \omterm{def:g5:series-parallel-circuits:parallel}{शाखा} का \omterm{ex:g3:first-electric-circuit:switch}{स्विच} खुला: \omterm{def:g3:first-electric-circuit:bulb}{बल्ब} 1 जलता है,
\omterm{def:g3:first-electric-circuit:bulb}{बल्ब} 2 अँधेरा।
\end{solution}

\begin{solution}{exo:g5:series-parallel-circuits:6}
\omterm{def:g4:series-circuits:series}{श्रेणी में} एक मरा हुआ \omterm{def:g3:first-electric-circuit:bulb}{बल्ब} (या बुझा दिया गया एक लैंप भी!) पूरे घर को
अँधेरा कर देता, और जितनी बत्तियाँ जुड़तीं सब उतनी ही मंद पड़ती जातीं।
समांतर जोड़ से मिलने वाली सुविधाएँ: हर कमरे का \omterm{ex:g3:first-electric-circuit:switch}{स्विच} सिर्फ़ अपने कमरे
पर राज करता है; रसोई का फुँका हुआ \omterm{def:g3:first-electric-circuit:bulb}{बल्ब} बाक़ी घर को जलता छोड़ देता है;
और हर लैंप पूरी चमक पर जलता है।
\end{solution}

\begin{solution}{exo:g5:series-parallel-circuits:7}
फूट से पहले --- यानी उस साझा सड़क पर जहाँ से धारा घर में घुसती है:
फ़्यूज़-पेटी का मुख्य \omterm{ex:g3:first-electric-circuit:switch}{स्विच}। वह पूरे घर की बिजली एक साथ काट देता है,
और मरम्मत तथा आपात स्थितियों में काम आता है।
\end{solution}

\begin{solution}{exo:g5:series-parallel-circuits:8}
\omterm{def:g5:series-parallel-circuits:parallel}{समांतर में}: छत की बत्ती बुझाने से सिरहाने का लैंप नहीं मरता, और लैंप
का प्लग निकालने से छत अँधेरी नहीं होती --- हर एक अपनी \omterm{def:g5:series-parallel-circuits:parallel}{शाखा} पर सवार
है।
\end{solution}

\begin{solution}{exo:g5:series-parallel-circuits:9}
हमारे सीधे-सादे स्विचों से डिज़ाइन बार-बार नाकाम होते हैं: \omterm{def:g3:first-electric-circuit:bulb}{बल्ब} के साथ
\omterm{def:g4:series-circuits:series}{श्रेणी में} लगा एक \omterm{ex:g3:first-electric-circuit:switch}{स्विच} सिर्फ़ एक सिरे से खोला जा सकता है; \omterm{def:g4:series-circuits:series}{श्रेणी में}
लगे दो सादे \omterm{ex:g3:first-electric-circuit:switch}{स्विच} \emph{दोनों} का बंद होना माँगते हैं (कोई भी सिरा
बत्ती बुझा सकता है, पर हमेशा जला नहीं सकता); और \omterm{def:g5:series-parallel-circuits:parallel}{समांतर में} लगे दो
\omterm{ex:g3:first-electric-circuit:switch}{स्विच} अँधेरे के लिए दोनों का खुला होना माँगते हैं। गलियारे वाली
सुविधा के लिए ऐसा \omterm{ex:g3:first-electric-circuit:switch}{स्विच} चाहिए जो सड़क तोड़ने के बजाय \emph{दो सड़कों
में से एक चुने} --- यानी दो-तरफ़ा \omterm{ex:g3:first-electric-circuit:switch}{स्विच}, जिसे असली गलियारे जोड़ियों
में काम में लेते हैं।
\end{solution}

\begin{solution}{exo:g5:series-parallel-circuits:10}
सुविधा: \omterm{def:g3:first-electric-circuit:bulb}{बल्ब} समांतर \omterm{def:g5:series-parallel-circuits:parallel}{शाखाओं} पर सवार हैं, इसलिए मरा हुआ \omterm{def:g3:first-electric-circuit:bulb}{बल्ब} सिर्फ़
ख़ुद को अँधेरा करता है। चेतावनी: हर ग़ायब \omterm{def:g3:first-electric-circuit:bulb}{बल्ब} एक \omterm{def:g5:series-parallel-circuits:parallel}{शाखा} हटा देता है,
और \omterm{def:g3:first-electric-circuit:battery}{बैटरी} (या लड़ी की छोटी बिजली-आपूर्ति) अपना पूरा ज़ोर बची हुई
\omterm{def:g5:series-parallel-circuits:parallel}{शाखाओं} में से भेजती है --- बहुत सारी शाखाएँ चली जाएँ तो बचे हुए हर
\omterm{def:g3:first-electric-circuit:bulb}{बल्ब} को उससे ज़्यादा मिलने लगता है जितने के लिए वह बना था, वह
ज़्यादा तपता है और जल्दी मर जाता है।
\end{solution}

\begin{solution}{exo:g5:series-parallel-circuits:11}
\omterm{def:g5:series-parallel-circuits:parallel}{शाखा} वाले \omterm{ex:g3:first-electric-circuit:switch}{स्विच} की अपनी सड़क पर: \omterm{def:g5:series-parallel-circuits:parallel}{शाखा} के \omterm{ex:g3:first-electric-circuit:switch}{स्विच} और \omterm{def:g3:first-electric-circuit:bulb}{बल्ब} 2 के बीच
(उनकी \omterm{def:g5:series-parallel-circuits:parallel}{शाखा} के बाक़ी परिपथ से मिलने से पहले)। वहाँ वह ठीक तभी जलता है
जब वह \omterm{def:g5:series-parallel-circuits:parallel}{शाखा} चल रही हो --- और वह \omterm{def:g3:first-electric-circuit:bulb}{बल्ब} 2 के साथ \omterm{def:g4:series-circuits:series}{श्रेणी में} बैठता है,
क्योंकि दोनों एक ही \omterm{def:g5:series-parallel-circuits:parallel}{शाखा} पर सवार हैं।
\end{solution}
